\chapter{Background and Related Work}
\label{ch:background}

This chapter surveys the four bodies of work this thesis draws on and
positions itself against: the Lean~4 ecosystem and its formal-library
practice; Petri-net and workflow-net soundness theory together with the
process-mining literature it feeds; the growing landscape of
LLM-assisted and automated theorem proving; and Rust verification and
extraction. It closes with an explicit statement of the gap between
these tracks that motivates this thesis's contribution.

\section{Lean 4 and Its Ecosystem}

Lean~4 itself, as a dependently typed theorem prover and programming
language with an elaboration and tactic framework distinct from Lean~3,
is described by de Moura and Ullrich~\cite{moura2021lean4}. The design
choice most relevant to this thesis is that Lean~4's kernel is small and
independently re-implementable: Lean4Lean, an independent kernel-in-Lean
typechecker, has been used to re-verify the reference kernel's behavior
and, in doing so, located and helped fix a soundness bug in the
reference implementation~\cite{lean4lean2024}. This matters for the
manufacturing pipeline described in this thesis because the pipeline's
own trust boundary is explicitly drawn \emph{above} the kernel: the
pipeline does not attempt to re-verify or extend kernel soundness, it
relies on the kernel as an already-scrutinized primitive and confines its
own apparatus (candidates, refusals, admissions, manifests) to the layer
of process around kernel calls.

Mathlib~\cite{mathlib2020} is the reference point for what a
large-scale, collaboratively maintained Lean formal library looks like in
practice: its quality regime is human review plus continuous
integration, applied uniformly across a very large and heterogeneous
contributor base. \textsf{procint} depends on Mathlib for its background
mathematics (finite types, order structures, and the like), but does not
attempt to replicate Mathlib's review-based quality regime; instead it
substitutes a mechanically checked admission gate for human review as the
primary standing-conferring step, a substitution this thesis argues is
appropriate specifically because \textsf{procint}'s corpus is
manufactured from a declared, machine-readable obligation graph rather
than authored freely.

CSLib~\cite{cslib2026} extends the Lean ecosystem toward computer
science proper --- labeled transition systems, traces, and bisimulation
among its concerns --- and \textsf{procint} builds directly on this layer
for parts of its semantics, in particular where process behavior is
described via transition and trace structures rather than purely via
algebraic Petri-net state equations. The presence of a computer-science
library alongside Mathlib's predominantly mathematical corpus is itself
evidence that Lean's applicability is already understood to extend past
pure mathematics; \textsf{procint} extends that applicability further,
into a domain --- process intelligence --- that neither Mathlib nor
CSLib targets directly.

\section{Petri Nets, Workflow-Net Soundness, and Process Mining}

The structural theory of Petri nets that \textsf{procint}'s
\texttt{ProcInt.Petri} namespace formalizes --- the state equation,
place and transition invariants, boundedness, and liveness --- follows
Murata's survey~\cite{murata1989}, still the standard reference
articulation of these properties and their relationships (for instance,
that liveness implies deadlock-freedom but not conversely).

Workflow nets, and the soundness property this thesis's crown-jewel
theorem addresses, are due to van der Aalst~\cite{vanderaalst1997}. The
specific result \textsf{procint} formalizes is van der Aalst's Lemma~8:
for a workflow net with a single source place, a single sink place, and
every node on some path from source to sink, soundness is equivalent to
liveness and boundedness of the net's short-circuited variant (the net
obtained by adding a transition from the sink back to the source). This
equivalence is exactly the theorem this thesis refers to throughout as
the crown-jewel theorem; Chapter~\ref{ch:procint} gives its Lean
statement and admission status in detail. The subsequent decidability
and classification landscape for workflow-net soundness --- which
soundness variants are decidable, and under what net restrictions --- is
surveyed in~\cite{vanderaalst2011soundness}, and the broader process
mining setting these soundness notions serve (discovery, conformance,
enhancement) is presented as a textbook treatment
in~\cite{vanderaalst2011pm}.

\textsf{procint}'s conformance-checking layer, which relates a process
model to observed event-log behavior via alignments, moves, and cost,
follows Carmona et al.'s treatment of conformance
checking~\cite{carmona2018}, which frames alignment-based techniques as
the technically preferred (if computationally heavier) alternative to
naive token-replay methods. Object-centric event data --- where a single
event may relate to several objects of possibly different types, rather
than to a single case identifier --- follows the OCEL~2.0
specification~\cite{ocel2}, and \textsf{procint}'s \texttt{ProcInt.Ocel}
namespace formalizes the entity-to-entity and object-to-object relations,
flattening, and lifecycle notions that specification defines.
Object-centric Petri nets, the net-level counterpart to object-centric
event logs, follow~\cite{ocpn2020}. Partially ordered process models,
formalized alongside process trees, DFGs, causal nets, BPMN, and Declare
in \textsf{procint}'s \texttt{ProcInt.Models} namespace, follow the POWL
language of Kourani and van Zelst~\cite{powl2023}.

To the best of a bounded search of the Lean, Coq, and Isabelle/AFP
library indexes conducted for this project (a negative search result,
not a proof of absence), none of workflow-net soundness, alignment-based
conformance, or OCEL had a prior public mechanization in a proof
assistant at the time \textsf{procint} was manufactured. A separate,
internal literature scan conducted for this project and recorded in this
repository's own research notes (\texttt{research-papers/README.md})
reaches a related but distinct observation: across the papers that scan
surveyed, process mining and formal (theorem-prover) verification have
historically progressed along near-parallel tracks, with comparatively
few works that bridge the two --- the scan's own phrasing is that the two
areas ``operate in parallel'' with essentially no bridge papers located.
That internal document is a project research note rather than a
peer-reviewed source, and this thesis treats its finding accordingly, as
a motivating internal observation rather than as a citable external
claim; but it is consistent with, and independently arrived at from, the
absence this thesis reports from the Lean/Coq/Isabelle index search
above.

\section{LLM-Assisted and Automated Theorem Proving}

A substantial and fast-moving literature addresses how language models
and automated provers can assist, rather than replace, interactive
theorem proving in systems like Lean. Two architectural patterns
recur. The first is in-process integration, where a language model is
called directly from within the tactic framework via a foreign-function
boundary, as in Lean Copilot~\cite{leancopilot2024}. The second is
external bridging to dedicated automated theorem provers or solvers
outside the interactive session, as in Lean-auto~\cite{leanauto2025},
which interfaces Lean~4 with external ATPs. \textsf{mfact} adopts neither
pattern as a trust mechanism: whatever automation --- language-model or
otherwise --- is used to produce a candidate declaration or proof
script, the candidate is untrusted data with respect to the pipeline's
own trust boundary, admitted only through the gate developed in
Chapter~\ref{ch:mfact}. This is a deliberate design choice rather than an
oversight of the in-process/external-bridge distinction: from the
pipeline's point of view, an in-process copilot suggestion and an
external solver's output are the same kind of object, a candidate, and
neither is privileged relative to a template-rendered or a
hand-authored one.

Search-based and learning-based automated proving more broadly is
represented by HOList~\cite{holist2019}, an environment for machine
learning applied to higher-order theorem proving, and by HyperTree Proof
Search~\cite{hypertree2022}, which frames proof search itself as a
tree-search problem amenable to neural guidance. TheoremLlama~\cite{theoremllama2024}
and APOLLO~\cite{apollo2025} represent, respectively, adapting a
general-purpose LLM into a Lean~4-specialized prover and orchestrating
collaboration between an LLM and Lean's own checking loop toward more
advanced formal reasoning. Infrastructure supporting these efforts
includes premise and declaration search over existing libraries, as in
LeanExplore~\cite{leanexplore2025}, and large synthetic theorem corpora
generated by exploring state-transition graphs, as
in~\cite{mathlib4dataset2025}. \textsf{procint} does not currently have
the problems this infrastructure addresses: its statement corpus is
curated from a declared RDF/Turtle catalog rather than mined from an
existing library or synthesized at scale, so premise search over a large
undifferentiated corpus and synthetic-theorem generation are, for this
project's current scope, orthogonal concerns rather than adjacent
solutions.

Two further strands bear directly on this thesis's central claim that
kernel admission alone is an insufficient standard for LLM-produced
formal content. Programming Language Confusion~\cite{langconfusion2025}
documents that code-generating LLMs can fail to keep target languages
straight even when producing superficially well-formed output, an
instance of the broader concern that syntactic plausibility and semantic
fidelity can diverge under automated generation. Work on understanding
and improving automated proof synthesis for interactive theorem
provers~\cite{tacticexpertalign2026} studies exactly the alignment
between what a synthesized tactic sequence proves and what a human
intended it to prove. Most directly, Meek et al.~\cite{meek2026formalizing}
formalize numerical analysis using a coding-agent pipeline and audit its
output for quality beyond kernel acceptance, documenting concretely that
compilation success can mask an incompletely stated multi-part theorem,
a silently added weakening hypothesis, or an unannounced parameter
restriction. Meek et al. is the closest prior work to \textsf{mfact} in
problem framing, and this thesis differs from it in mechanism rather
than diagnosis, as elaborated in the Research Gap section below.

\section{Rust Verification and Extraction}

The correspondence rail developed in Chapter~\ref{ch:correspondence}
extracts a bounded Rust implementation core into a Lean-facing semantic
image via Aeneas~\cite{ho2022aeneas}, a tool for Rust verification by
functional translation, whose Charon front end~\cite{protzenko2024charon}
separates Rust syntax recovery from the semantic translation step
proper. Aeneas is designed to target whole Rust programs under a full
memory model; this thesis's use of it is narrower, applied to a single,
dependency-free arithmetic module, and --- consistent with the
untrusted-producer discipline argued for throughout this thesis ---
extraction's output is treated as a candidate Lean statement, never as a
proof: standing for a correspondence claim still requires kernel
admission of a hand-authored correspondence theorem stated over the
extracted declaration, not merely a same-shaped statement over
independently hand-written types. Chapter~\ref{ch:correspondence}
returns to this point in more mechanical detail, including the guard
that blocks a correspondence obligation from reporting proven status
while its extraction binding remains incomplete.

\section{Research Gap}

Taken individually, none of the tracks surveyed above addresses the
problem this thesis targets. The Lean ecosystem (Mathlib, CSLib,
Lean~4Lean) supplies library scale, computer-science-relevant semantics
infrastructure, and an independently checked kernel, but its standard
quality regime is human review, not a mechanically checked admission
pipeline applied uniformly to agent-, template-, and human-produced
candidates alike. The Petri-net, workflow-net, and process-mining
literature supplies the mathematical canon --- soundness, alignments,
OCEL --- that this thesis's first manufactured vertical formalizes, but
that literature is, per both this project's own bounded library-index
search and its internal research scan, essentially disjoint from formal
verification practice: process mining and theorem proving have
progressed as near-parallel tracks with few bridge works, so no existing
mechanization supplies workflow-net soundness or OCEL as a Lean~4 target
to build on or compare against directly. The LLM-assisted and automated
theorem-proving literature supplies both the capability this thesis takes
as its starting condition (agents and models can produce plausible formal
candidates) and, in Meek et al.~\cite{meek2026formalizing} particularly, a
sharp diagnosis of the resulting risk, but that diagnosis is delivered as
a post-hoc semantic audit layered over agent output after the fact, not
as a typed, structurally enforced pipeline that a candidate must pass
through before a standing claim can be recorded at all. The Rust
verification and extraction literature (Aeneas, Charon) supplies a
mechanism for turning implementation code into a proof-facing semantic
object, but does not by itself address how the resulting extracted
statement's standing, or the standing of any hand-authored correspondence
theorem built on it, should be tracked, gated, or receipted over time.

The gap this thesis addresses sits at the intersection of these tracks:
a domain formalization effort (process intelligence) that (a) has, to
this project's search, no prior proof-assistant mechanization to inherit
conventions from; (b) is produced with substantial LLM and agent
assistance, and therefore inherits the standing problem Meek et al.
diagnose; and (c) is manufactured, rather than hand-authored, from a
declared obligation graph, so that the standing problem must be solved
structurally --- as a typed admission pipeline with typed refusals,
manifests, and replayable receipts --- rather than remedied afterward by
a separate audit pass. Where Meek et al.'s audit is a quality assessment
performed \emph{over} already-produced agent output, this thesis's
contribution, \textsf{mfact}, is a pipeline in which candidates,
refusals, and admissions are themselves Lean data types, so that the
proven/stated distinction is enforced by a release gate rather than
estimated after the fact. The two approaches are, as
paper/main.tex's Related Work section observes, complementary rather
than competing: the residual risk that a kernel-admitted statement is
nonetheless a faithless or vacuous rendering of its informal target is
exactly the risk a semantic audit like Meek et al.'s is suited to catch,
and exactly the risk this thesis's pipeline, by design, does not claim
to eliminate. Chapter~\ref{ch:discussion} returns to this residual risk
directly, rather than treating kernel admission as a substitute for it.
